\documentclass[11pt,a4paper]{article}

% -------------------------------
% Encoding & fonts (pdflatex-safe)
% -------------------------------
\usepackage[utf8]{inputenc}
\usepackage[T1]{fontenc}
\usepackage{lmodern}

% -------------------------------
% Layout
% -------------------------------
\usepackage[
a4paper,
left=0.3in,
right=0.3in,
top=0.75in,
bottom=1in
]{geometry}

% -------------------------------------
% TikZ
% -------------------------------------
\usepackage{tikz}
\usetikzlibrary{
	positioning,
	calc,
	arrows.meta,
	decorations.pathreplacing,
	shapes.geometric
}

% -------------------------------
% Tables, lists, colors
% -------------------------------
\usepackage{array}
\usepackage{booktabs}
\usepackage{makecell}
\usepackage{multicol}
\usepackage{enumitem}
\usepackage{xcolor}

% -------------------------------
% Math
% -------------------------------
\usepackage{amsmath,amssymb}

% -------------------------------
% Header & Footer
% -------------------------------
\usepackage{fancyhdr}
\usepackage{lastpage}

\setlength{\headheight}{28pt}

\pagestyle{fancy}
\fancyhf{}

\fancyhead[L]{\textbf{Lebanese University\\ULFG II - 1h30}}
\fancyhead[C]{\textbf{Concurrent Programming\\Midterm Exam}}
\fancyhead[R]{\textbf{Date: 29/04/2026\\Closed Book Exam}}

\fancyfoot[L]{Dr M. AOUDE}
\fancyfoot[C]{}
\fancyfoot[R]{\thepage/\pageref{LastPage}}

% -------------------------------
% Helpers & macros
% -------------------------------
\newcounter{qnum}

\newcommand{\Q}{%
	\stepcounter{qnum}%
	\noindent\textbf{Q\theqnum. }%
}

\newcommand{\chkraw}[1]{%
	\tikz[baseline=(b.center)]{
		\node[
		draw,
		minimum width=#1,
		minimum height=#1,
		inner sep=0pt
		] (b) {};
	}%
}

\newcommand{\ansrow}[1]{%
	#1 & \chkraw{4mm} & \chkraw{4mm} & \chkraw{4mm} & \chkraw{4mm} \\ \hline
}

\newcommand{\key}[1]{\underline{#1}}

% -------------------------------
% Listings: Java / Concurrent Programming Style
% -------------------------------
\usepackage{listings}

\definecolor{codebg}{RGB}{245,245,245}
\definecolor{codeborder}{RGB}{210,210,210}

\lstdefinestyle{java}{
	language=Java,
	backgroundcolor=\color{codebg},
	basicstyle=\ttfamily\small,
	frame=single,
	rulecolor=\color{codeborder},
	showstringspaces=false,
	breaklines=true,
	keywordstyle=\color{blue!70!black}\bfseries,
	commentstyle=\itshape\color{gray!70!black},
	stringstyle=\color{brown!70!black},
	columns=fullflexible,
	tabsize=2
}

\lstset{style=java}

% Safe inline code
\DeclareRobustCommand{\code}[1]{\texttt{\detokenize{#1}}}

% -------------------------------
% MCQ Choices
% -------------------------------
\newcommand{\choices}[4]{%
	\setlength{\tabcolsep}{4pt}%
	\renewcommand{\arraystretch}{1.15}%
	\begin{tabular}{|p{0.46\linewidth}|p{0.46\linewidth}|}
		\hline
		A) #1 & B) #2 \\ \hline
		C) #3 & D) #4 \\ \hline
	\end{tabular}%
}

\begin{document}
	% -------------------------------
	% QCM Answer Sheet
	% -------------------------------
	\begin{center}
		\Large \textbf{QCM Answer Sheet (1–55)}\\[4pt]
		\small Mark your answer with \textbf{X} inside the chosen box.
	\end{center}
	
	\noindent
	\begin{minipage}[t]{0.49\linewidth}
		\renewcommand{\arraystretch}{1.7}
		\begin{tabular}{|>{\centering\arraybackslash}m{1.3cm}
				|>{\centering\arraybackslash}m{1.3cm}
				|>{\centering\arraybackslash}m{1.3cm}
				|>{\centering\arraybackslash}m{1.3cm}
				|>{\centering\arraybackslash}m{1.3cm}|}
			\hline
			\textbf{\#} & \textbf{A} & \textbf{B} & \textbf{C} & \textbf{D} \\ \hline
			\ansrow{1}\ansrow{2}\ansrow{3}\ansrow{4}\ansrow{5}
			\ansrow{6}\ansrow{7}\ansrow{8}\ansrow{9}\ansrow{10}
			\ansrow{11}\ansrow{12}\ansrow{13}\ansrow{14}\ansrow{15}
			\ansrow{16}\ansrow{17}\ansrow{18}\ansrow{19}\ansrow{20}
			\ansrow{21}\ansrow{22}\ansrow{23}\ansrow{24}\ansrow{25}
			
		\end{tabular}
	\end{minipage}
	\hfill
	\begin{minipage}[t]{0.49\linewidth}
		\renewcommand{\arraystretch}{1.7}
		\begin{tabular}{|>{\centering\arraybackslash}m{1.3cm}
				|>{\centering\arraybackslash}m{1.3cm}
				|>{\centering\arraybackslash}m{1.3cm}
				|>{\centering\arraybackslash}m{1.3cm}
				|>{\centering\arraybackslash}m{1.3cm}|}
			\hline
			\textbf{\#} & \textbf{A} & \textbf{B} & \textbf{C} & \textbf{D} \\ \hline
			\ansrow{26}\ansrow{27}\ansrow{28}\ansrow{29}\ansrow{30}
			\ansrow{31}\ansrow{32}\ansrow{33}
			\ansrow{34}\ansrow{35}\ansrow{36}\ansrow{37}\ansrow{38}
			\ansrow{39}\ansrow{40}\ansrow{41}\ansrow{42}\ansrow{43}
			\ansrow{44}\ansrow{45}\ansrow{46}\ansrow{47}\ansrow{48}\ansrow{49}
			\ansrow{50}
		\end{tabular}
	\end{minipage}
	
	\newpage
	\begin{center}
		{\Large \textbf{Exercise I  -- MCQ Quiz (55 Questions)}}\\[2mm]
	\end{center}
	
	\vspace{2mm}

\Q In a single-threaded server, what happens when one request is being processed? \textbf{(Que se passe-t-il dans un serveur mono-thread lorsqu’une requête est en cours de traitement ?) }

\choices
{Other clients are immediately processed in parallel \textbf{(Les autres clients sont immédiatement traités en parallèle)}}
{New clients must wait until the current request finishes \textbf{(Les nouveaux clients doivent attendre la fin de la requête en cours)}}
{The server automatically creates new threads \textbf{(Le serveur crée automatiquement de nouveaux threads)}}
{The operating system rejects all new clients \textbf{(Le système d’exploitation rejette tous les nouveaux clients)}}

\Q What is the main weakness of this loop? \textbf{(Quelle est la principale faiblesse de cette boucle ?) }

\begin{lstlisting}[language=Java]
	while (true) {
		Socket client = serverSocket.accept();
		handleRequest(client);
	}
\end{lstlisting}

\choices
{It processes only one request at a time \textbf{(Elle traite une seule requête à la fois)}}
{It cannot accept sockets \textbf{(Elle ne peut pas accepter des sockets)}}
{It requires a database \textbf{(Elle nécessite une base de données)}}
{It always crashes under one client \textbf{(Elle plante toujours avec un seul client)}}

\Q When 50 clients connect to a single-threaded server, the main problem is usually: \textbf{(Lorsque 50 clients se connectent à un serveur mono-thread, quel est le principal problème ?) }

\choices
{Lower memory usage \textbf{(Utilisation mémoire plus faible)}}
{Queue growth and increasing latency \textbf{(Croissance de la file d’attente et augmentation de la latence)}}
{No CPU usage \textbf{(Aucune utilisation du CPU)}}
{Faster response time \textbf{(Temps de réponse plus rapide)}}

\Q The expression \code{throughput = requests / second} measures: \textbf{(Que mesure l’expression débit = requêtes / seconde ?) }

\choices
{How many requests are completed per second \textbf{(Le nombre de requêtes traitées par seconde)}}
{How many threads exist \textbf{(Le nombre de threads existants)}}
{How many files are open \textbf{(Le nombre de fichiers ouverts)}}
{How much RAM is installed \textbf{(La quantité de mémoire RAM installée)}}

\Q For I/O-bound workloads, increasing the number of threads can help because: \textbf{(Pour les charges I/O, pourquoi augmenter le nombre de threads peut aider ?) }

\choices
{Threads can make progress while others are waiting for I/O \textbf{(Les threads peuvent progresser pendant que d’autres attendent les entrées/sorties)}}
{CPU instructions become mathematically free \textbf{(Les instructions CPU deviennent gratuites mathématiquement)}}
{Locks disappear automatically \textbf{(Les verrous disparaissent automatiquement)}}
{Network latency becomes zero \textbf{(La latence réseau devient nulle)}}

\Q What is the CPU-bound trap? \textbf{(Quel est le piège des tâches liées au CPU ?) }

\choices
{Adding more threads always gives unlimited speedup \textbf{(Ajouter plus de threads donne toujours un gain illimité)}}
{Adding too many threads to CPU-heavy work may not improve performance \textbf{(Ajouter trop de threads pour des tâches CPU peut ne pas améliorer les performances)}}
{CPU-bound programs never use the CPU \textbf{(Les programmes CPU n’utilisent jamais le CPU)}}
{Thread pools cannot run Java code \textbf{(Les pools de threads ne peuvent pas exécuter du code Java)}}

\Q The engineering formula for thread pool tuning depends strongly on: \textbf{(De quoi dépend fortement le réglage d’un pool de threads ?) }

\choices
{The balance between waiting time and compute time \textbf{(L’équilibre entre temps d’attente et temps de calcul)}}
{The color of the terminal \textbf{(La couleur du terminal)}}
{The file name of the Java class \textbf{(Le nom du fichier Java)}}
{The number of comments in the code \textbf{(Le nombre de commentaires dans le code)}}

\Q Amdahl's Law explains: \textbf{(Que décrit la loi d’Amdahl ?) }

\choices
{The maximum possible speedup when only part of a program is parallelizable \textbf{(Le gain maximal possible lorsqu’une partie seulement est parallélisable)}}
{How to create a socket \textbf{(Comment créer un socket)}}
{How to write SQL queries \textbf{(Comment écrire des requêtes SQL)}}
{How to disable threads \textbf{(Comment désactiver les threads)}}

\Q If a program has a large sequential part, then even infinite cores will: \textbf{(Si une grande partie est séquentielle, que se passe-t-il même avec des cœurs infinis ?) }

\choices
{Give unlimited speedup \textbf{(Donner un gain illimité)}}
{Still be limited by the sequential part \textbf{(Rester limité par la partie séquentielle)}}
{Remove all latency \textbf{(Supprimer toute latence)}}
{Make synchronization unnecessary \textbf{(Rendre la synchronisation inutile)}}

\Q In Little's Law, \(L\) represents: \textbf{(Dans la loi de Little, que représente L ?) }

\choices
{Average number of items in the system \textbf{(Nombre moyen d’éléments dans le système)}}
{Number of CPU cores \textbf{(Nombre de cœurs CPU)}}
{Length of the source code \textbf{(Longueur du code source)}}
{Number of Java classes \textbf{(Nombre de classes Java)}}

\Q Context switching cost refers to: \textbf{(À quoi correspond le coût de changement de contexte ?) }

\choices
{The overhead of switching the CPU from one thread to another \textbf{(Le coût du passage du CPU d’un thread à un autre)}}
{The cost of buying a new computer \textbf{(Le coût d’achat d’un ordinateur)}}
{The time needed to compile comments \textbf{(Le temps pour compiler les commentaires)}}
{The size of the database \textbf{(La taille de la base de données)}}

\Q What does p50 latency represent? \textbf{(Que représente la latence p50 ?) }

\choices
{The median response time \textbf{(Le temps de réponse médian)}}
{The maximum possible response time \textbf{(Le temps de réponse maximal)}}
{The number of threads \textbf{(Le nombre de threads)}}
{The number of failed requests only \textbf{(Le nombre de requêtes échouées uniquement)}}
	\Q What does p95 latency help reveal? \textbf{(Que révèle la latence p95 ?) }
	
	\choices
	{Slow responses experienced by the worst 5 percent of requests \textbf{(Les réponses lentes subies par les 5\% pires requêtes)}}
	{Only the fastest request \textbf{(Seulement la requête la plus rapide)}}
	{The number of CPU cores \textbf{(Le nombre de cœurs CPU)}}
	{The Java version \textbf{(La version de Java)}}
	
	\Q A critical section is: \textbf{(Qu’est-ce qu’une section critique ?) }
	
	\choices
	{Code that accesses shared data and must be protected \textbf{(Code qui accède à des données partagées et doit être protégé)}}
	{Code that is never executed \textbf{(Code qui n’est jamais exécuté)}}
	{A section used only for printing titles \textbf{(Une section utilisée uniquement pour afficher des titres)}}
	{A Java package name \textbf{(Un nom de package Java)}}
	
	\Q In Java, how can you create a thread? \textbf{(Comment peut-on créer un thread en Java ?) }
	
	\choices
	{By extending Thread or implementing Runnable \textbf{(En étendant Thread ou en implémentant Runnable)}}
	{By creating a Socket \textbf{(En créant un socket)}}
	{By calling System.exit() \textbf{(En appelant System.exit())}}
	{By using SQL queries \textbf{(En utilisant des requêtes SQL)}}
	
	\Q What happens when \code{start()} is called on a thread? \textbf{(Que se passe-t-il lorsqu’on appelle start() sur un thread ?) }
	
	\choices
	{A new thread begins execution and calls run() \textbf{(Un nouveau thread commence l’exécution et appelle run())}}
	{The program terminates \textbf{(Le programme se termine)}}
	{The thread is deleted \textbf{(Le thread est supprimé)}}
	{The CPU is stopped \textbf{(Le CPU s’arrête)}}
	
	\Q Calling \code{run()} directly instead of \code{start()}: \textbf{(Appeler run() directement au lieu de start() :) }
	
	\choices
	{Executes in the current thread (no concurrency) \textbf{(S’exécute dans le thread courant (pas de concurrence))}}
	{Creates multiple threads \textbf{(Crée plusieurs threads)}}
	{Causes a deadlock \textbf{(Provoque un interblocage)}}
	{Improves parallelism \textbf{(Améliore le parallélisme)}}
	
	\Q A race condition occurs when: \textbf{(Quand une condition de course se produit-elle ?) }
	
	\choices
	{Multiple threads access shared data without proper synchronization \textbf{(Plusieurs threads accèdent à des données partagées sans synchronisation adéquate)}}
	{Only one thread executes \textbf{(Un seul thread s’exécute)}}
	{A program runs too fast \textbf{(Un programme s’exécute trop rapidement)}}
	{A variable is constant \textbf{(Une variable est constante)}}
	
	\Q Consider:
	\begin{lstlisting}[language=Java]
		counter++;
	\end{lstlisting}
	
	Why is this unsafe in multithreading? \textbf{(Pourquoi cela est-il dangereux en multithreading ?) }
	
	\choices
	{It is not atomic and can be interrupted \textbf{(Ce n’est pas une opération atomique et elle peut être interrompue)}}
	{It is too slow \textbf{(C’est trop lent)}}
	{It uses too much memory \textbf{(Cela utilise trop de mémoire)}}
	{It cannot compile \textbf{(Cela ne compile pas)}}
	
	\Q Race conditions are dangerous because: \textbf{(Pourquoi les conditions de course sont-elles dangereuses ?) }
	
	\choices
	{They produce unpredictable results \textbf{(Elles produisent des résultats imprévisibles)}}
	{They always crash immediately \textbf{(Elles provoquent toujours un crash immédiat)}}
	{They prevent compilation \textbf{(Elles empêchent la compilation)}}
	{They remove threads \textbf{(Elles suppriment les threads)}}
	
	\Q Mutual exclusion ensures that: \textbf{(Que garantit l’exclusion mutuelle ?) }
	
	\choices
	{Only one thread accesses a critical section at a time \textbf{(Un seul thread accède à une section critique à la fois)}}
	{All threads run simultaneously \textbf{(Tous les threads s’exécutent simultanément)}}
	{No thread can run \textbf{(Aucun thread ne peut s’exécuter)}}
	{Threads do not exist \textbf{(Les threads n’existent pas)}}
	
	\Q In Java, which keyword provides mutual exclusion? \textbf{(Quel mot-clé en Java fournit l’exclusion mutuelle ?) }
	
	\choices
	{\code{synchronized} \textbf{(synchronized)}}
	{\code{public} \textbf{(public)}}
	{\code{static} \textbf{(static)}}
	{\code{return} \textbf{(return)}}
	
\Q What is the main purpose of a lock? \textbf{(Quel est le but principal d’un verrou ?) }

\choices
{To prevent simultaneous access to shared resources \textbf{(Empêcher l’accès simultané aux ressources partagées)}}
{To speed up printing \textbf{(Accélérer l’affichage)}}
{To reduce memory usage \textbf{(Réduire l’utilisation mémoire)}}
{To stop threads permanently \textbf{(Arrêter définitivement les threads)}}

\Q A deadlock occurs when: \textbf{(Quand un interblocage se produit-il ?) }

\choices
{Threads wait forever for resources held by each other \textbf{(Les threads attendent indéfiniment des ressources détenues par d’autres)}}
{Threads execute too fast \textbf{(Les threads s’exécutent trop rapidement)}}
{Memory is full \textbf{(La mémoire est pleine)}}
{A program ends \textbf{(Un programme se termine)}}

\Q Which condition is required for deadlock? \textbf{(Quelle condition est nécessaire pour un interblocage ?) }

\choices
{Circular waiting \textbf{(Attente circulaire)}}
{No threads \textbf{(Aucun thread)}}
{Single variable \textbf{(Une seule variable)}}
{Fast CPU \textbf{(CPU rapide)}}

\Q Example of deadlock: \textbf{(Exemple d’interblocage :) }

\choices
{Thread A holds Lock1 and waits for Lock2, while Thread B holds Lock2 and waits for Lock1 \textbf{(Le thread A détient Lock1 et attend Lock2, tandis que le thread B détient Lock2 et attend Lock1)}}
{A thread prints a message \textbf{(Un thread affiche un message)}}
{A loop increments a variable \textbf{(Une boucle incrémente une variable)}}
{A file is read \textbf{(Un fichier est lu)}}

\Q One way to avoid deadlock is: \textbf{(Une façon d’éviter un interblocage est :) }

\choices
{Always acquire locks in the same order \textbf{(Toujours acquérir les verrous dans le même ordre)}}
{Remove all threads \textbf{(Supprimer tous les threads)}}
{Increase memory \textbf{(Augmenter la mémoire)}}
{Use more print statements \textbf{(Utiliser plus d’instructions d’affichage)}}

\Q Deadlock results in: \textbf{(Un interblocage entraîne :) }

\choices
{System freeze or blocked execution \textbf{(Blocage du système ou exécution figée)}}
{Faster execution \textbf{(Exécution plus rapide)}}
{Lower CPU usage only \textbf{(Seulement une baisse de l’utilisation CPU)}}
{Compilation failure \textbf{(Échec de compilation)}}

\Q What is visibility in concurrency? \textbf{(Qu’est-ce que la visibilité en concurrence ?) }

\choices
{Whether one thread sees updates made by another thread \textbf{(Si un thread voit les mises à jour faites par un autre thread)}}
{How fast a thread runs \textbf{(La vitesse d’exécution d’un thread)}}
{How many threads exist \textbf{(Le nombre de threads existants)}}
{The size of memory \textbf{(La taille de la mémoire)}}

\Q Without proper synchronization: \textbf{(Sans synchronisation adéquate :) }

\choices
{Threads may see stale values \textbf{(Les threads peuvent voir des valeurs obsolètes)}}
{Threads always see latest values \textbf{(Les threads voient toujours les dernières valeurs)}}
{Threads stop working \textbf{(Les threads cessent de fonctionner)}}
{Program compiles faster \textbf{(Le programme compile plus rapidement)}}

\Q Which keyword helps ensure visibility in Java? \textbf{(Quel mot-clé assure la visibilité en Java ?) }

\choices
{\code{volatile} \textbf{(volatile)}}
{\code{int} \textbf{(int)}}
{\code{class} \textbf{(class)}}
{\code{main} \textbf{(main)}}

\Q The Java Memory Model defines: \textbf{(Que définit le modèle mémoire Java ?) }

\choices
{Rules for how threads interact through memory \textbf{(Les règles d’interaction des threads via la mémoire)}}
{Database schemas \textbf{(Les schémas de bases de données)}}
{File system structure \textbf{(La structure du système de fichiers)}}
{Network protocols \textbf{(Les protocoles réseau)}}

\Q Happens-before relationship ensures: \textbf{(La relation happens-before garantit :) }

\choices
{Correct ordering and visibility of operations between threads \textbf{(Un ordre correct et la visibilité des opérations entre threads)}}
{Faster compilation \textbf{(Une compilation plus rapide)}}
{Lower memory usage \textbf{(Une utilisation mémoire réduite)}}
{Thread deletion \textbf{(La suppression des threads)}}

\Q The Producer-Consumer problem involves: \textbf{(Le problème producteur-consommateur implique :) }

\choices
{Producers generating data and consumers using it \textbf{(Des producteurs générant des données et des consommateurs les utilisant)}}
{Only producers writing files \textbf{(Seulement des producteurs écrivant des fichiers)}}
{Only consumers reading memory \textbf{(Seulement des consommateurs lisant la mémoire)}}
{Threads that do not interact \textbf{(Des threads qui n’interagissent pas)}}

\Q What structure is typically used between producer and consumer? \textbf{(Quelle structure est utilisée entre producteur et consommateur ?) }

\choices
{Shared buffer (queue) \textbf{(Tampon partagé (file d’attente))}}
{File system \textbf{(Système de fichiers)}}
{Stack trace \textbf{(Trace de pile)}}
{Compiler \textbf{(Compilateur)}}

\Q If the buffer is full, the producer should: \textbf{(Si le tampon est plein, le producteur doit :) }

\choices
{Wait \textbf{(Attendre)}}
{Overwrite data \textbf{(Écraser les données)}}
{Terminate immediately \textbf{(Se terminer immédiatement)}}
{Ignore the buffer \textbf{(Ignorer le tampon)}}

\Q If the buffer is empty, the consumer should: \textbf{(Si le tampon est vide, le consommateur doit :) }

\choices
{Wait \textbf{(Attendre)}}
{Create data \textbf{(Créer des données)}}
{Delete the buffer \textbf{(Supprimer le tampon)}}
{Restart the program \textbf{(Redémarrer le programme)}}
	\Q The main challenge in this problem is: \textbf{(Quel est le principal défi dans ce problème ?) }
	
	\choices
	{Synchronization between threads \textbf{(La synchronisation entre les threads)}}
	{Memory allocation only \textbf{(L’allocation mémoire uniquement)}}
	{Compilation errors \textbf{(Les erreurs de compilation)}}
	{File reading speed \textbf{(La vitesse de lecture des fichiers)}}
	
	\Q Condition variables are typically used with: \textbf{(Les variables de condition sont généralement utilisées avec :) }
	
	\choices
	{Locks / mutexes \textbf{(Des verrous / mutex)}}
	{Arrays \textbf{(Des tableaux)}}
	{Files \textbf{(Des fichiers)}}
	{Compilers \textbf{(Des compilateurs)}}
	
	\Q What problem do condition variables help solve? \textbf{(Quel problème les variables de condition aident-elles à résoudre ?) }
	
	\choices
	{Busy waiting inefficiency \textbf{(L’inefficacité de l’attente active (busy waiting))}}
	{Syntax errors \textbf{(Les erreurs de syntaxe)}}
	{Memory leaks \textbf{(Les fuites mémoire)}}
	{Compilation time \textbf{(Le temps de compilation)}}
	
	\Q One necessary condition for deadlock is: \textbf{(Quelle est une condition nécessaire pour un interblocage ?) }
	
	\choices
	{Mutual exclusion \textbf{(Exclusion mutuelle)}}
	{No threads \textbf{(Aucun thread)}}
	{Infinite memory \textbf{(Mémoire infinie)}}
	{No locks \textbf{(Aucun verrou)}}
	
	\Q Which is a deadlock prevention strategy? \textbf{(Quelle est une stratégie de prévention des interblocages ?) }
	
	\choices
	{Lock ordering \textbf{(Ordonnancement des verrous)}}
	{Increasing CPU speed \textbf{(Augmenter la vitesse du CPU)}}
	{Reducing code size \textbf{(Réduire la taille du code)}}
	{Removing variables \textbf{(Supprimer des variables)}}
	
	\Q Breaking circular wait prevents: \textbf{(Rompre l’attente circulaire permet d’éviter :) }
	
	\choices
	{Deadlock \textbf{(Les interblocages)}}
	{Race conditions \textbf{(Les conditions de course)}}
	{Compilation errors \textbf{(Les erreurs de compilation)}}
	{Thread creation \textbf{(La création de threads)}}
	
	\Q Resource allocation graphs are used to: \textbf{(Les graphes d’allocation de ressources sont utilisés pour :) }
	
	\choices
	{Analyze deadlocks \textbf{(Analyser les interblocages)}}
	{Compile programs \textbf{(Compiler des programmes)}}
	{Store variables \textbf{(Stocker des variables)}}
	{Manage files \textbf{(Gérer des fichiers)}}
\Q Deadlock avoidance differs from prevention because: \textbf{(En quoi l’évitement de l’interblocage diffère-t-il de la prévention ?) }

\choices
{It dynamically checks safe states \textbf{(Il vérifie dynamiquement les états sûrs)}}
{It removes all locks \textbf{(Il supprime tous les verrous)}}
{It disables threads \textbf{(Il désactive les threads)}}
{It ignores resources \textbf{(Il ignore les ressources)}}

\Q Starvation occurs when: \textbf{(Quand la famine (starvation) se produit-elle ?) }

\choices
{A thread never gets access to resources \textbf{(Un thread n’obtient jamais l’accès aux ressources)}}
{All threads finish quickly \textbf{(Tous les threads se terminent rapidement)}}
{Memory is full \textbf{(La mémoire est pleine)}}
{CPU is idle \textbf{(Le CPU est inactif)}}

\Q Fairness in scheduling means: \textbf{(Que signifie l’équité dans l’ordonnancement ?) }

\choices
{All threads eventually get a chance to execute \textbf{(Tous les threads finissent par avoir une chance de s’exécuter)}}
{Only one thread runs \textbf{(Un seul thread s’exécute)}}
{Threads are removed \textbf{(Les threads sont supprimés)}}
{Execution is random \textbf{(L’exécution est aléatoire)}}

\Q Which mechanism helps reduce starvation? \textbf{(Quel mécanisme aide à réduire la famine ?) }

\choices
{Fair locks / scheduling policies \textbf{(Verrous équitables / politiques d’ordonnancement équitables)}}
{Infinite loops \textbf{(Boucles infinies)}}
{Removing threads \textbf{(Suppression des threads)}}
{Increasing file size \textbf{(Augmenter la taille des fichiers)}}

\Q Priority scheduling may cause: \textbf{(L’ordonnancement par priorité peut provoquer :) }

\choices
{Starvation of low-priority threads \textbf{(La famine des threads de faible priorité)}}
{Compilation errors \textbf{(Des erreurs de compilation)}}
{Memory leaks \textbf{(Des fuites mémoire)}}
{No execution \textbf{(Aucune exécution)}}

\Q Aging technique is used to: \textbf{(La technique de vieillissement est utilisée pour :) }

\choices
{Increase priority of waiting threads over time \textbf{(Augmenter la priorité des threads en attente avec le temps)}}
{Delete old threads \textbf{(Supprimer les anciens threads)}}
{Reduce memory \textbf{(Réduire la mémoire)}}
{Stop execution \textbf{(Arrêter l’exécution)}}
	
	% ============================================================
	\section*{Exercise II -- Code Tracing and Completion}
	% ============================================================
	
	\textbf{Code Tracing: Race Condition}
	
	Consider the following Java code:
	
	\begin{lstlisting}[language=Java]
		class Counter {
			private int count = 0;
			
			public void increment() {
				count++;
			}
			
			public int getCount() {
				return count;
			}
		}
		
		public class TestRace {
			public static void main(String[] args) throws Exception {
				Counter c = new Counter();
				
				Thread t1 = new Thread(() -> {
					for (int i = 0; i < 1000; i++) {
						c.increment();
					}
				});
				
				Thread t2 = new Thread(() -> {
					for (int i = 0; i < 1000; i++) {
						c.increment();
					}
				});
				
				t1.start();
				t2.start();
				
				t1.join();
				t2.join();
				
				System.out.println(c.getCount());
			}
		}
	\end{lstlisting}
	
\textbf{Question II :}  
What is the expected behavior of this program? Should the output always be \code{2000}? Explain briefly.
\textbf{(Quel est le comportement attendu de ce programme ? La sortie doit-elle toujours être \code{2000} ? Expliquez brièvement.)}\vspace{12pt}
	
	
	\noindent\makebox[\linewidth]{\dotfill}\par\vspace{6pt}
	\vspace{0.5cm}
	\noindent\makebox[\linewidth]{\dotfill}\par\vspace{6pt}
	\vspace{0.5cm}
	\noindent\makebox[\linewidth]{\dotfill}\par\vspace{6pt}
	\vspace{0.5cm}
	\noindent\makebox[\linewidth]{\dotfill}\par\vspace{6pt}
	\vspace{0.5cm}
	\noindent\makebox[\linewidth]{\dotfill}\par\vspace{6pt}
	\vspace{0.5cm}
	\noindent\makebox[\linewidth]{\dotfill}\par\vspace{6pt}
	\vspace{0.5cm}
	\noindent\makebox[\linewidth]{\dotfill}\par\vspace{6pt}
	\vspace{0.5cm}
	\noindent\makebox[\linewidth]{\dotfill}\par
	\vspace{0.5cm}
	
	% ============================================================
	\textbf{Question III:}  
% ============================================================
\textbf{Question III:}  
\textbf{Code Completion: Thread-Safe Bank Account}
\textbf{(Complétion de code : Compte bancaire thread-safe)}

We want the following behavior:
\textbf{(On souhaite le comportement suivant :)}

\begin{itemize}
	\item Two threads perform deposits on the same bank account.
	\textbf{(Deux threads effectuent des dépôts sur le même compte bancaire.)}
	
	\item The final balance must always be correct.
	\textbf{(Le solde final doit toujours être correct.)}
	
	\item No race condition should occur.
	\textbf{(Aucune condition de course ne doit se produire.)}
\end{itemize}

Complete the missing part of the code.
\textbf{(Complétez la partie manquante du code.)}

\begin{lstlisting}[language=Java]
	class BankAccount {
		private int balance = 0;
		
		public void deposit(int amount) {
			// TODO: add the missing synchronization code here
			
			
			
			
			
			
			
			
			
			
			
			
			
			
			
			
			
			
			balance += amount;
		}
		
		public int getBalance() {
			return balance;
		}
	}
\end{lstlisting}

\textbf{Required:}  
Rewrite only the \code{deposit()} method so that the account becomes thread-safe.  
\textbf{(Réécrivez uniquement la méthode \code{deposit()} pour rendre le compte thread-safe.)}

\vspace{3cm}
\end{document}